% GENERATED FILE. Edit canonical lessons, then run npm run generate:books.
% canonical-source-hash: fnv1a64:e761b2d31574c64f
% canonical-lessons: GU-C508-lokhand, GU-C508-luhar, GU-C508-bhaththi, GU-C508-tapavvu, GU-C508-eran, GU-R508-review-501-508

\chapter{The Blacksmith's Fire}
\label{ch:gu-the-blacksmith-s-fire}
\hlchaptermodality{508}

\begin{chapteropening}
\textbf{By the end of this chapter:} \emph{I can talk about a blacksmith, his fire and his anvil.}

Complete the last lesson of chapter 508: I can talk about a blacksmith, his fire and his anvil.
\end{chapteropening}

\section[\texorpdfstring{\gu{લોખંડ} (lokhanḍ)}{લોખંડ (lokhanḍ)}]{\gu{લોખંડ} (lokhanḍ) — iron}

\label{lesson:GU-C508-lokhand}



\begin{quote}
Before the new one: say the Gujarati for a patch (on cloth), then the Gujarati for embroidery.
\end{quote}

\subsection*{You'll want to know: \gu{લોખંડ}}
\textbf{\gu{લોખંડ}} — \emph{lokhanḍ} — \textquotedblleft{}iron\textquotedblright{}.

\begin{quote}
\gu{આ} \gu{દરવાજો} \gu{લોખંડનો} \gu{છે}.
\end{quote}

\begin{quote}
\emph{ā darvājo lokhanḍno che.}
\end{quote}

\begin{quote}
This gate is made of iron.
\end{quote}

A neuter noun.

\begin{cousinweb}
\textbf{\gu{લોખંડ}} (\emph{lokhanḍ}) comes from \emph{*lōhōpaskara}, \textquotedblleft{}iron tools\textquotedblright{}, a form CDIAL builds from \emph{lōhá}, \textquotedblleft{}iron\textquotedblright{}, and \emph{upaskara}, \textquotedblleft{}tool\textquotedblright{}; CDIAL lists the Gujarati word as \textquotedblleft{}tools, iron, ironware\textquotedblright{} (CDIAL 11171).
\end{cousinweb}

\begin{grammarlens}[title={What you've built}]
One more word for describing what you lived and felt.
\end{grammarlens}

\subsection*{Your turn}
\begin{itemize}
\raggedright
  \item \emph{Say it:} \emph{lokhanḍ}
  \item \emph{Say it:} \emph{lokhanḍ}, once more
  \item \emph{Say it:} \emph{ā darvājo lokhanḍno che.} — This gate is made of iron.
  \item \emph{Recall:} say the Gujarati for a patch (on cloth), then the Gujarati for embroidery, then say \emph{lokhanḍ} again
\end{itemize}

\begin{tcolorbox}[breakable,colback=teal!4,colframe=teal!35!black,title={Before you move on}]
What does \gu{લોખંડ} mean? (\textquotedblleft{}Iron\textquotedblright{}.) Say it once more.
\end{tcolorbox}
\section[\texorpdfstring{\gu{લુહાર} (luhār)}{લુહાર (luhār)}]{\gu{લુહાર} (luhār) — a blacksmith}

\label{lesson:GU-C508-luhar}



\begin{quote}
Before the new one: say the Gujarati for embroidery, then the Gujarati for iron.
\end{quote}

\subsection*{You'll want to know: \gu{લુહાર}}
\textbf{\gu{લુહાર}} — \emph{luhār} — \textquotedblleft{}a blacksmith\textquotedblright{}.

\begin{quote}
\gu{લુહાર} \gu{લોખંડમાંથી} \gu{ઓજારો} \gu{બનાવે} \gu{છે}.
\end{quote}

\begin{quote}
\emph{luhār lokhanḍmā̃thī ojāro banāve che.}
\end{quote}

\begin{quote}
The blacksmith makes tools out of iron.
\end{quote}

A masculine noun.

\begin{cousinweb}
\textbf{\gu{લુહાર}} (\emph{luhār}) goes back to Sanskrit \emph{lōhakāra}, \textquotedblleft{}an iron-worker\textquotedblright{}, built from \emph{lōha}, \textquotedblleft{}iron\textquotedblright{}, and \emph{kāra}, \textquotedblleft{}maker\textquotedblright{}; CDIAL gives the Gujarati form as \emph{lavār} (CDIAL 11159).
\end{cousinweb}

\begin{grammarlens}[title={What you've built}]
One more word for describing what you lived and felt.
\end{grammarlens}

\subsection*{Your turn}
\begin{itemize}
\raggedright
  \item \emph{Say it:} \emph{luhār}
  \item \emph{Say it:} \emph{luhār}, once more
  \item \emph{Say it:} \emph{luhār lokhanḍmā̃thī ojāro banāve che.} — The blacksmith makes tools out of iron.
  \item \emph{Recall:} say the Gujarati for embroidery, then the Gujarati for iron, then say \emph{luhār} again
\end{itemize}

\begin{tcolorbox}[breakable,colback=teal!4,colframe=teal!35!black,title={Before you move on}]
What does \gu{લુહાર} mean? (\textquotedblleft{}A blacksmith\textquotedblright{}.) Say it once more.
\end{tcolorbox}
\section[\texorpdfstring{\gu{ભઠ્ઠી} (bhaṭhṭhī)}{ભઠ્ઠી (bhaṭhṭhī)}]{\gu{ભઠ્ઠી} (bhaṭhṭhī) — a furnace}

\label{lesson:GU-C508-bhaththi}



\begin{quote}
Before the new one: say the Gujarati for iron, then the Gujarati for a blacksmith.
\end{quote}

\subsection*{You'll want to know: \gu{ભઠ્ઠી}}
\textbf{\gu{ભઠ્ઠી}} — \emph{bhaṭhṭhī} — \textquotedblleft{}a furnace\textquotedblright{}.

\begin{quote}
\gu{લુહારની} \gu{ભઠ્ઠીમાં} \gu{આગ} \gu{સળગે} \gu{છે}.
\end{quote}

\begin{quote}
\emph{luhārnī bhaṭhṭhīmā̃ āg saḷge che.}
\end{quote}

\begin{quote}
A fire burns in the blacksmith's furnace.
\end{quote}

A feminine noun.

\begin{cousinweb}
\textbf{\gu{ભઠ્ઠી}} (\emph{bhaṭhṭhī}) is a furnace or a kiln. CDIAL lists Panjabi \emph{bhaṭṭhī}, \textquotedblleft{}a furnace\textquotedblright{}, and Hindi \emph{bhaṭṭhā}, \textquotedblleft{}a kiln\textquotedblright{}, under Sanskrit \emph{bhráṣṭra}, \textquotedblleft{}a frying pan, a gridiron\textquotedblright{} (CDIAL 9656); it gives no Gujarati form for this word, so its exact history is not settled.
\end{cousinweb}

\begin{grammarlens}[title={What you've built}]
One more word for describing what you lived and felt.
\end{grammarlens}

\subsection*{Your turn}
\begin{itemize}
\raggedright
  \item \emph{Say it:} \emph{bhaṭhṭhī}
  \item \emph{Say it:} \emph{bhaṭhṭhī}, once more
  \item \emph{Say it:} \emph{luhārnī bhaṭhṭhīmā̃ āg saḷge che.} — A fire burns in the blacksmith's furnace.
  \item \emph{Recall:} say the Gujarati for iron, then the Gujarati for a blacksmith, then say \emph{bhaṭhṭhī} again
\end{itemize}

\begin{tcolorbox}[breakable,colback=teal!4,colframe=teal!35!black,title={Before you move on}]
What does \gu{ભઠ્ઠી} mean? (\textquotedblleft{}A furnace\textquotedblright{}.) Say it once more.
\end{tcolorbox}
\section[\texorpdfstring{\gu{તપાવવું} (tapāvvũ)}{તપાવવું (tapāvvũ)}]{\gu{તપાવવું} (tapāvvũ) — to heat (something)}

\label{lesson:GU-C508-tapavvu}



\begin{quote}
Before the new one: say the Gujarati for a blacksmith, then the Gujarati for a furnace.
\end{quote}

\subsection*{You'll want to know: \gu{તપાવવું}}
\textbf{\gu{તપાવવું}} — \emph{tapāvvũ} — \textquotedblleft{}to heat (something)\textquotedblright{}.

\begin{quote}
\gu{લુહારે} \gu{ભઠ્ઠીમાં} \gu{લોખંડ} \gu{તપાવ્યું}.
\end{quote}

\begin{quote}
\emph{luhāre bhaṭhṭhīmā̃ lokhanḍ tapāvyũ.}
\end{quote}

\begin{quote}
The blacksmith heated the iron in the furnace.
\end{quote}

Past, with \gu{તેણે} (he or she …): \gu{તેણે} \gu{લોખંડ} \gu{તપાવ્યું}, \emph{teṇe lokhanḍ tapāvyũ}.

\begin{cousinweb}
\textbf{\gu{તપાવવું}} (\emph{tapāvvũ}) is built from \textbf{\gu{તપવું}}, \textquotedblleft{}to be hot\textquotedblright{}, with the causative \emph{-āv-}: to make something hot. CDIAL files the family under Sanskrit \emph{tapyati}, \textquotedblleft{}is hot\textquotedblright{} (CDIAL 5684).
\end{cousinweb}

\begin{grammarlens}[title={What you've built}]
One more word for describing what you lived and felt.
\end{grammarlens}

\subsection*{Your turn}
\begin{itemize}
\raggedright
  \item \emph{Say it:} \emph{tapāvvũ}
  \item \emph{Say it:} \emph{tapāvvũ}, once more
  \item \emph{Say it:} \emph{luhāre bhaṭhṭhīmā̃ lokhanḍ tapāvyũ.} — The blacksmith heated the iron in the furnace.
  \item \emph{Recall:} say the Gujarati for a blacksmith, then the Gujarati for a furnace, then say \emph{tapāvvũ} again
\end{itemize}

\begin{tcolorbox}[breakable,colback=teal!4,colframe=teal!35!black,title={Before you move on}]
What does \gu{તપાવવું} mean? (\textquotedblleft{}To heat (something)\textquotedblright{}.) Say it once more.
\end{tcolorbox}
\section[\texorpdfstring{\gu{એરણ} (eraṇ)}{એરણ (eraṇ)}]{\gu{એરણ} (eraṇ) — an anvil}

\label{lesson:GU-C508-eran}



\begin{quote}
Before the new one: say the Gujarati for a furnace, then the Gujarati for to heat (something).
\end{quote}

\subsection*{You'll want to know: \gu{એરણ}}
\textbf{\gu{એરણ}} — \emph{eraṇ} — \textquotedblleft{}an anvil\textquotedblright{}.

\begin{quote}
\gu{લુહાર} \gu{એરણ} \gu{પર} \gu{ગરમ} \gu{લોખંડ} \gu{મૂકે} \gu{છે}.
\end{quote}

\begin{quote}
\emph{luhār eraṇ par garam lokhanḍ mūke che.}
\end{quote}

\begin{quote}
The blacksmith puts the hot iron on the anvil.
\end{quote}

A feminine noun.

\begin{cousinweb}
\textbf{\gu{એરણ}} (\emph{eraṇ}) is inherited from Sanskrit \emph{adhikaraṇī}, \textquotedblleft{}an anvil\textquotedblright{}, and CDIAL lists the Gujarati word (CDIAL 252).
\end{cousinweb}

\begin{grammarlens}[title={What you've built}]
One more word for describing what you lived and felt.
\end{grammarlens}

\subsection*{Your turn}
\begin{itemize}
\raggedright
  \item \emph{Say it:} \emph{eraṇ}
  \item \emph{Say it:} \emph{eraṇ}, once more
  \item \emph{Say it:} \emph{luhār eraṇ par garam lokhanḍ mūke che.} — The blacksmith puts the hot iron on the anvil.
  \item \emph{Recall:} say the Gujarati for a furnace, then the Gujarati for to heat (something), then say \emph{eraṇ} again
\end{itemize}

\begin{tcolorbox}[breakable,colback=teal!4,colframe=teal!35!black,title={Before you move on}]
What does \gu{એરણ} mean? (\textquotedblleft{}An anvil\textquotedblright{}.) Say it once more.
\end{tcolorbox}
\section[(dialogue)]{Looking back: from The Craftsman's Bench to The Blacksmith's Fire}

\label{lesson:GU-R508-review-501-508}



\begin{quote}
Say the Gujarati for wood and a craftsman.
\end{quote}

\subsection*{Your turn}
Say these aloud:

\begin{itemize}
\raggedright
  \item wood, a craftsman, a tool of a trade, to know how, a profession
  \item a saw, to hollow out, a hammer, a nail, to hammer in
  \item a carpenter's plane, to smooth, glue, a crack, to split (wood)
  \item a potter, a potter's wheel, mud plaster, to knead clay, a clay lamp
  \item the cotton plant, raw cotton, a spinning wheel, to spin (thread), khadi
  \item a weaver, a strand, a weave, to stretch tight, to wind round
  \item to thread (a needle), a pearl, to mend (a tear), a patch (on cloth), embroidery
  \item iron, a blacksmith, a furnace, to heat (something), an anvil
\end{itemize}

\begin{tcolorbox}[breakable,colback=teal!4,colframe=teal!35!black,title={Before you move on}]
Wood? (\textbf{\gu{લાકડું}}, \emph{lākḍũ}.) A craftsman? (\textbf{\gu{કારીગર}}, \emph{kārīgar}.) A tool of a trade? (\textbf{\gu{ઓજાર}}, \emph{ojār}.) To know how? (\textbf{\gu{આવડવું}}, \emph{āvaḍvũ}.) A profession? (\textbf{\gu{વ્યવસાય}}, \emph{vyavsāy}.) A saw? (\textbf{\gu{કરવત}}, \emph{karvat}.) To hollow out? (\textbf{\gu{કોરવું}}, \emph{korvũ}.) A hammer? (\textbf{\gu{હથોડી}}, \emph{hathoḍī}.) A nail? (\textbf{\gu{મેખ}}, \emph{mekh}.) To hammer in? (\textbf{\gu{ઠોકવું}}, \emph{ṭhokvũ}.) A carpenter's plane? (\textbf{\gu{રંદો}}, \emph{rando}.) To smooth? (\textbf{\gu{મઠારવું}}, \emph{maṭhārvũ}.) Glue? (\textbf{\gu{ગુંદર}}, \emph{gundar}.) A crack? (\textbf{\gu{તિરાડ}}, \emph{tirāḍ}.) To split (wood)? (\textbf{\gu{ચીરવું}}, \emph{cīrvũ}.) A potter? (\textbf{\gu{કુંભાર}}, \emph{kumbhār}.) A potter's wheel? (\textbf{\gu{ચાકડો}}, \emph{cākḍo}.) Mud plaster? (\textbf{\gu{ગારો}}, \emph{gāro}.) To knead clay? (\textbf{\gu{ગૂંદવું}}, \emph{gū̃dvũ}.) A clay lamp? (\textbf{\gu{કોડિયું}}, \emph{koḍiyũ}.) The cotton plant? (\textbf{\gu{કપાસ}}, \emph{kapās}.) Raw cotton? (\textbf{\gu{રૂ}}, \emph{rū}.) A spinning wheel? (\textbf{\gu{રેંટિયો}}, \emph{rẽṭiyo}.) To spin (thread)? (\textbf{\gu{કાંતવું}}, \emph{kā̃tvũ}.) Khadi? (\textbf{\gu{ખાદી}}, \emph{khādī}.) A weaver? (\textbf{\gu{વણકર}}, \emph{vaṇkar}.) A strand? (\textbf{\gu{તાંતણો}}, \emph{tā̃tṇo}.) A weave? (\textbf{\gu{વણાટ}}, \emph{vaṇāṭ}.) To stretch tight? (\textbf{\gu{તાણવું}}, \emph{tāṇvũ}.) To wind round? (\textbf{\gu{વીંટવું}}, \emph{vī̃ṭvũ}.) To thread (a needle)? (\textbf{\gu{પરોવવું}}, \emph{parovvũ}.) A pearl? (\textbf{\gu{મોતી}}, \emph{motī}.) To mend (a tear)? (\textbf{\gu{સાંધવું}}, \emph{sā̃dhvũ}.) A patch (on cloth)? (\textbf{\gu{થીંગડું}}, \emph{thī̃gḍũ}.) Embroidery? (\textbf{\gu{ભરતકામ}}, \emph{bharatkām}.) Iron? (\textbf{\gu{લોખંડ}}, \emph{lokhanḍ}.) A blacksmith? (\textbf{\gu{લુહાર}}, \emph{luhār}.) A furnace? (\textbf{\gu{ભઠ્ઠી}}, \emph{bhaṭhṭhī}.) To heat (something)? (\textbf{\gu{તપાવવું}}, \emph{tapāvvũ}.) An anvil? (\textbf{\gu{એરણ}}, \emph{eraṇ}.)
\end{tcolorbox}
